% ----------
% Tech report template
% 
% ----------
\documentclass[10pt,a4paper]{article}
\usepackage{geometry}
%\usepackage[utf8x]{inputenc}
\usepackage[T1]{fontenc}
\usepackage[english]{babel}
\usepackage[runin]{abstract}
\usepackage{titling}
\usepackage{listings}
\usepackage{booktabs}
\usepackage{mlmodern}
%\usepackage{piton} % Without this enable, it might not work for the listing package. Switch to LuaLateX if you want to use this option, however. 
\usepackage{fancyhdr}
%\usepackage{helvet}
\usepackage{csquotes}
\usepackage{graphicx}
\usepackage{blindtext}
\usepackage{parskip}
\usepackage{etoolbox}
%\usepackage[scaled=0.90]{helvet} % Font 1
%\usepackage[scaled=0.85]{beramono} % Font 2
\usepackage{xcolor}
\usepackage{hyperref}
% Bibliography
\usepackage[backend=biber,style=alphabetic]{biblatex}
\addbibresource{bibb.bib} %Imports bibliography file

\input{preamble.tex}

% ----------
% Variables
% ----------

\title{\textbf{Laboratory Formation or Directive Proposal Template}, v1  \\ {\small Introduction and specification of \textbf{structural specification and topic classification}}} % Full title of your tech report
\runningtitle{Reformation documentations} % Short title
\author{Fujimiya Amane \\ {\small Laboratory Principal Director and Researcher}} % Full list of authors
\runningauthor{Amane et al.} % Short list of authors
\affiliation{} % Affiliation e.g. University or Company
\department{Your Department} % Department or Office
\memoid{General-Sub-Serial-Additional Code} % ID of the tech report
\theyear{Year} % year of the tech report
\mydate{Full month, year} %the date
\usepackage{tcolorbox}
\usepackage{fancyvrb}

\tcbuselibrary{minted,breakable,xparse,skins}

\definecolor{bg}{gray}{0.95}
%\DeclareTCBListing{mintedbox}{O{}m!O{}}{%
%  breakable=true,
%  listing engine=minted,
%  listing only,
%  minted language=#2,
%  minted style=default,
%  minted options={%
%    linenos,
%    gobble=0,
%    breaklines=true,
%    breakafter=,,
%    fontsize=\small,
%    numbersep=8pt,
%    #1},
%  boxsep=0pt,
%  left skip=0pt,
%  right skip=0pt,
%  left=25pt,
%  right=0pt,
%  top=3pt,
%  bottom=3pt,
%  arc=5pt,
%  leftrule=0pt,
%  rightrule=0pt,
%  bottomrule=2pt,
%  toprule=2pt,
%  colback=bg,
%  colframe=orange!70,
%  enhanced,
%  overlay={%
%    \begin{tcbclipinterior}
%    \fill[orange!20!white] (frame.south west) rectangle ([xshift=20pt]frame.north west);
%    \end{tcbclipinterior}},
%  #3}

% Python preset environment. 
\newcommand{\pyobject}[1]{\fbox{\color{red}{\texttt{#1}}}}

\NewDocumentCommand{\codeword}{v}{%
\texttt{\textcolor{blue}{#1}}%
}
% Inline python highlighting (Might not look great, but eh)
\lstset{language=Python,keywordstyle={\bfseries \color{blue}}}


% ----------
% actual document
% ----------
\begin{document}
\maketitle

\begin{abstract}
    \noindent
    This is a template for all RHINELAB-RINALAB \textbf{directive} and \textbf{structural specification}, changes, mandate, or laboratory \textit{formation statements}
    
    This is \textbf{version 1} of the template. Subjected to change. 

    \keywords{Keyword1, Keyword2, Keyword3}
\end{abstract}

\vspace{2.5cm}



\thispagestyle{firstpage}

\pagebreak

% ----------
% End of first page
% ----------

\newgeometry{} % Redefine geometries (normal margins)

\section{Specification}
\label{sec:spec}

Provide specification details, metaprogramming. 

\subsection{Document code, priority}

Provide document codes, additional document code, purpose, organization, authorship and anyone contributing to the report, identification of purpose (proposal, or research, indexing 01). Priority code, additional conditions. Does it override anything, or add to anything?

\subsection{Body of issue}

Who issue, team or departmental attachment. Since this is a \textbf{structural document}, write down of which director or high-ranking member is responsible for this proposal template.  

\subsection{Date and Version}

Version specification across versions, and similarly date of construction. \textbf{Do not remove old dates}, for versioning record. 

\section{Introduction}

This is a specific template that is \textit{only used} for and by members up to the role of Director and above, without immediate effect or overriding priority. The purpose of it is simple - to provide specification on either: 
\begin{enumerate}
    \item Laboratory components or functionality formation proposal or of sort. 
    \item Specification on structural and changes of components of the laboratory. 
    \item Directive, i.e. \textbf{direction and structural function classification} of the laboratory or its components, the way moving forward or else given of topic, strategy, targets and aims of the components. 
    \item Miscellaneous structural proposal on events, planning or any given \textit{dynamical activities} (that is something definitely not planned in the overall structures and only related to stuffs that happens on the fly). 
\end{enumerate}
Those would be what the template can allow you to \textit{proposal something}. It is directly distinct from the proposal template, of which is there for the majority of the time, reserved and resorted to research, academic activities, and projects. 

For the remainder of the introduction section, you are to be expected to fill in the following information.

\subsection{Responsibility}

Who would be the responsible personnel, scholar, individuals or components and segments of the laboratory in responsible for such function and changes outlined here? Who would be the people in drafting of the documents? 

\subsection{Risk assessment}

How much would this proposal affects the general structural of the laboratory? What is the added \textbf{structural formation and components} to accommodate the directive and the idea thereof? What is the capacities and resource topology to be used to support? How successful or rather of the finishing availability you are seeing from the proposal? 

\section{Central proposal structure}

The following contents are required in drafting and specification of the document. We note that the majority of information required are in response to the checkup of \textit{realization and laboratory assessment}, that is, if you know about the intricacy and real situation that is surrounded of the laboratory, or else, and your navigational aptitude in response to such. 

\subsection{Nature and current situation}

Because of the setup for this proposal is required to have its \textit{object of modification or proposed mutation}, the setting of which necessitate such changes, or, under certainty or uncertainty thereof, be given of the vision inquired for the laboratory or its \textbf{splinter components} related to the laboratory (like collaborations, or sub-lab or semi-dependent entities), the following will need to be answered: 

\subsubsection{Object and Identifier}

What is the structures, component, interest, individuals, groups, research components or any \textbf{subject of interest} that this proposal is aiming at? 

\subsubsection{Situation}

What is the situation, general context, \textbf{states and current form} of the subject specified above? 

\subsubsection{Necessity}

By then, what is the perceived \textbf{necessity} thereof that gave rise to the propensity of this proposal? What is the parts \textit{of the laboratory interest} that makes it to be conceded to consider the proposal on its own ground? 

\subsubsection{Recheck, v1}

Are you aware of the \textbf{philosophical tradition and commitment} of the laboratory? If so, validate and recheck of your proposal to the perception and predisposition thereof, that it acquires necessary validation and compatibility given such, of the philosophical pillars. 

Are you aware of the culture and code of conduct? If so, recheck on that specifically. 

\subsection{Structural details}



\clearpage

\section*{Acknowledgements (optional)}

Thanks again to \textbf{Person 1} and \textbf{Affiliation A} for their financial support. Thanks \textbf{Person 2}, \textbf{Person 3}, etc. for their participation in drafting this document. Further information. 

% ----------
% Bibliography
% ----------

% \bibliography{../biblio.bib}
% \bibliographystyle{abbrv}

\appendix

\section{An appendix}

Appendix to more information. Analysis, cost process, etc. 

\section{Another appendix}

\subsection*{Subsectioning in appendix}

Some more text, and a list:

\blindenumerate



\section{Citations}
Items that are cited: \textit{The \LaTeX\ Companion} book \cite{latexcompanion} together with Einstein's journal paper \cite{einstein} and Dirac's book \cite{dirac}---which are physics-related items. Next, citing two of Knuth's books: \textit{Fundamental Algorithms} \cite{knuth-fa} and \textit{The Art of Computer Programming} \cite{knuthwebsite}.


Remember to change this to your respective bibliography. Usually, the file that holds the bibliography in the path is \texttt{bibb.bib}, so beware of that. The below part where there is \texttt{printbibliography} does it work - it prints the bibliography. Cite them using the \texttt{cite} key seen in the above citation. 

\medskip

\printbibliography %Prints bibliography


\end{document}

% ----------
